\section*{Chapter \ref{ch:b2:measurement-statistics} --- Statistics of Chemical Measurement}

\begin{solution}{exo:b2:measurement-statistics:1}
$\bar V = 62.41/5 = \qty{12.482}{mL}$; $\sum(V_i - \bar V)^2 = 0.00328$, $s = \sqrt{0.00328/4} \approx
\qty{0.029}{mL}$; $s/\sqrt5 \approx \qty{0.013}{mL}$.
\end{solution}

\begin{solution}{exo:b2:measurement-statistics:2}
$t = 2.776$ for 4 degrees of freedom: $12.482 \pm 2.776 \times 0.0128$, that is $(12.48 \pm 0.04)$~mL.
\end{solution}

\begin{solution}{exo:b2:measurement-statistics:3}
The relative uncertainties of the mass and of the volume are 0.04~\% and 0.06~\%; combined in
quadrature,
\[\frac{u(c)}{c} = \sqrt{\Bigl(\frac{0.0002}{0.5000}\Bigr)^2 + \Bigl(\frac{0.15}{250.0}\Bigr)^2}
\approx 0.07~\%.\]
\end{solution}

\begin{solution}{exo:b2:measurement-statistics:4}
$\bar x = 1.5$, $\bar y = 0.31$, $S_{xx} = 5$, $S_{xy} = 0.99$: $b = 0.198$, $a = 0.31 - 0.198 \times 1.5
= 0.013$.
\end{solution}

\begin{solution}{exo:b2:measurement-statistics:5}
$|10.12 - 10.00|/(0.08/\sqrt5) = 0.12/0.036 \approx 3.4 > 2.776$: the difference is significant, the
method is biased (by about \qty{+0.12}{mg/L}).
\end{solution}

\begin{solution}{exo:b2:measurement-statistics:6}
$\mathrm{pH} = -\log_{10}c$, $\mathrm d\,\mathrm{pH}/\mathrm dc = -1/(c\ln10)$:
$u(\mathrm{pH}) = (u(c)/c)/\ln10 = 0.02/2.303 \approx 0.009$.
\end{solution}

\begin{solution}{exo:b2:measurement-statistics:7}
$x_{\mathrm{LOD}} = 3 \times 0.0012/0.0098 \approx \qty{0.37}{\micro g/L}$; $x_{\mathrm{LOQ}} = 10
\times 0.0012/0.0098 \approx \qty{1.2}{\micro g/L}$.
\end{solution}

\begin{solution}{exo:b2:measurement-statistics:8}
The line gives \qty{3.11}{mg/L} in the measured solutions; the sample had been diluted fivefold:
$3.11 \times 50.0/10.0 \approx \qty{15.6}{mg/L}$.
\end{solution}

\begin{solution}{exo:b2:measurement-statistics:9}
\qty{0.05}{mg/L} is the \omterm{def:b2:measurement-statistics:validation}{repeatability}, \qty{0.15}{mg/L} the \omterm{def:b2:measurement-statistics:validation}{reproducibility}. Each laboratory has its
own small \omterm{def:b2:measurement-statistics:validation}{bias} (calibration, standards, instrument); these biases differ between laboratories and add
to the scatter, so \omterm{def:b2:measurement-statistics:validation}{reproducibility} is always the larger.
\end{solution}

\begin{solution}{exo:b2:measurement-statistics:10}
$\partial S/\partial a = 0$ reads $\sum(y_i - a - bx_i) = 0$: the \omterm{def:b2:measurement-statistics:calibration}{residuals} sum to zero. Dividing by
$n$: $\bar y = a + b\bar x$, so $(\bar x, \bar y)$ is on the line.
\end{solution}

\begin{solution}{exo:b2:measurement-statistics:11}
One step:
\[\sqrt{(0.006/1.000)^2 + (0.08/100.0)^2} \approx 0.61~\%.\]
One tenfold step:
\[\sqrt{(0.02/10.00)^2 + (0.08/100.0)^2} \approx 0.22~\%,\]
and two of them in quadrature give $0.22\sqrt2 \approx 0.30~\%$. The two-step route is twice as precise: the small pipette dominates the one-step
route.
\end{solution}

\begin{solution}{exo:b2:measurement-statistics:12}
$z = 1.6/\sqrt{0.4^2 + 0.5^2} \approx 2.5 > 2$: not compatible; one laboratory (at least) has a \omterm{def:b2:measurement-statistics:validation}{bias}.
With expanded uncertainties ($k = 2$), the first gives $9.6 \pm 0.8$, which includes 10; the second
$11.2 \pm 1.0$, entirely above 10. Before any verdict on the water, the disagreement must be resolved,
for instance with a reference material.
% ledger: who:lead
\end{solution}

\begin{solution}{pb:b2:measurement-statistics:1}
\textbf{1.} The acid changes the atomisation and the response; standards and samples must have the same
matrix for the slope to apply.
\textbf{2.} $\bar x = 10$, $\bar y = 0.1004$, $S_{xx} = 250$, $S_{xy} = 2.445$: $b = 0.00978$ per
\unit{\micro g/L}, $a = 0.1004 - 0.0978 = 0.0026$.
\textbf{3.} $+0.0004$, $-0.0005$, $+0.0006$, $-0.0013$, $+0.0008$.
\textbf{4.} Signs alternate, no curvature: the line is adequate over 0--\qty{20}{\micro g/L}.
\textbf{5.} $s_{y/x} = \sqrt{3.1 \times 10^{-6}/3} \approx 0.0010$.
\textbf{6.} The blank (acid, water, furnace) gives a small absorbance of its own.
\textbf{7.} The sensitivity: absorbance per \unit{\micro g/L} of lead.
\textbf{8.} $\bar y_0 = 0.1010$; $s = 0.0030$.
\textbf{9.} $x_0 = (0.1010 - 0.0026)/0.00978 \approx \qty{10.06}{\micro g/L}$.
\textbf{10.} $s_{x_0} = (0.0010/0.00978)\sqrt{1/3 + 1/5 + (0.1010 - 0.1004)^2/(0.00978^2 \times 250)}
\approx 0.102 \times 0.730 \approx \qty{0.076}{\micro g/L}$.
\textbf{11.} $n - 2 = 3$; $t = 3.182$.
\textbf{12.} $10.06 \pm 3.182 \times 0.076 = (10.06 \pm 0.24)$~\unit{\micro g/L}.
\textbf{13.} The term $1/m$ falls from 1 to $1/3$: $s_{x_0}$ drops by about a third.
\textbf{14.} $f = 50.50/50.0 = 1.010$; $10.06 \times 1.010 \approx \qty{10.16}{\micro g/L}$.
\textbf{15.} $f = 1 + V_2/V_1$: $u^2(f) = (u(V_2)/V_1)^2 + (V_2u(V_1)/V_1^2)^2 = (0.005/50.0)^2 +
(0.50 \times 0.03/2500)^2$, $u(f) \approx 1.0 \times 10^{-4}$, a relative 0.01~\%.
\textbf{16.} No: 0.01~\% against $0.076/10.06 \approx 0.75~\%$.
\textbf{17.} $s_{\mathrm{blank}} \approx 0.00115$; $x_{\mathrm{LOD}} = 3 \times 0.00115/0.00978 \approx
\qty{0.35}{\micro g/L}$.
\textbf{18.} $10 \times 0.00115/0.00978 \approx \qty{1.2}{\micro g/L}$.
\textbf{19.} Yes, by far.
\textbf{20.} Yes: from 9.92 to \qty{10.40}{\micro g/L} in the tap water.
\textbf{21.} $(0.104 - 0.0026)/0.00978 \times 1.010 \approx \qty{10.5}{\micro g/L}$, and, read alone,
a confident ``above the guideline'' that the data do not support.
\textbf{22.} The rule used to build the interval captures the true value in 95~\% of the series to
which it is applied; it is not a probability about this one value.
\textbf{23.} More replicate readings and more standards near \qty{10}{\micro g/L} (the terms $1/m$ and
$1/n$), or a more sensitive method; the interval shrinks only slowly, as $1/\sqrt{\text{number}}$.
\textbf{24.} By analysing a certified reference water, or by spiking the sample with a known amount of
lead and checking the recovery.
\textbf{25.} Near \qty{10}{\micro g/L} the uncertainty of routine methods is a few per cent, as here: a
lower guideline value could not be checked reliably by ordinary laboratories.
\textbf{26.} $\boldsymbol{(10.16 \pm 0.24)}$~\unit{\micro g/L} at 95~\%: the interval contains
\qty{10}{\micro g/L}, so the measurement neither shows the water to exceed the guideline nor to meet it
with a margin.
\end{solution}
